\section{总结与延伸}

\subsection{讲者的核心总结}

Manning 在课程结尾总结了本讲的核心线索：

\begin{enumerate}
    \item 人类语言是最复杂的符号系统，也是 AI 最具挑战性的领域之一
    \item 传统的词义表示（WordNet、one-hot）无法有效捕捉语义相似性
    \item 分布式语义假说提供了全新的视角：``词义 = 上下文''
    \item Word2Vec 通过简单的``预测上下文''任务，从大量文本中学习出高质量的词向量
    \item 训练过程的核心是梯度下降——通过微积分指导参数调整方向
\end{enumerate}

\subsection{全课知识图谱}

\begin{center}
\begin{tikzpicture}[node distance=0.8cm, every node/.style={draw, rounded corners, minimum width=3.5cm, minimum height=0.7cm, font=\small}]
\node (lang) {人类语言与词义};
\node (trad) [below=of lang] {传统表示方法};
\node (dist) [below=of trad] {分布式语义假说};
\node (w2v) [below=of dist] {Word2Vec 算法};
\node (opt) [below=of w2v] {梯度下降优化};
\draw[->, thick] (lang) -- (trad);
\draw[->, thick] (trad) -- (dist);
\draw[->, thick] (dist) -- (w2v);
\draw[->, thick] (w2v) -- (opt);
\node[draw=none, right=1cm of lang, text width=5.5cm, font=\scriptsize] {交流、思维支架、知识载体};
\node[draw=none, right=1cm of trad, text width=5.5cm, font=\scriptsize] {WordNet（手工、不完整）、One-hot（正交、无相似性）};
\node[draw=none, right=1cm of dist, text width=5.5cm, font=\scriptsize] {Firth: ``You shall know a word by the company it keeps''};
\node[draw=none, right=1cm of w2v, text width=5.5cm, font=\scriptsize] {Skip-gram、Softmax 概率、中心词/上下文词};
\node[draw=none, right=1cm of opt, text width=5.5cm, font=\scriptsize] {负对数似然、链式法则、SGD};
\end{tikzpicture}
\end{center}

\subsection{关键 Takeaways}

\begin{importantbox}{五条核心要点}
\begin{enumerate}
    \item \textbf{语言是思维的支架}：人类语言不仅用于交流，更是高层次认知的基础工具
    \item \textbf{分布式语义是词表示的正确范式}：词的含义由其上下文决定，而非人工标注的关系
    \item \textbf{词向量 = 可计算的语义}：稠密低维向量使得语义相似性可以通过点积/余弦相似度直接度量
    \item \textbf{Softmax 是万能的概率化工具}：将任意实数向量转化为概率分布，在深度学习中无处不在
    \item \textbf{``观察值 $-$ 期望值'' 驱动学习}：Word2Vec 的梯度具有直观的统计含义，通过不断缩小预测与观察之间的差距来学习
\end{enumerate}
\end{importantbox}

\subsection{拓展阅读}

\begin{itemize}
    \item Mikolov et al., 2013. \textit{Efficient Estimation of Word Representations in Vector Space}: \url{https://arxiv.org/abs/1301.3781} —— Word2Vec 原始论文
    \item Mikolov et al., 2013. \textit{Distributed Representations of Words and Phrases and their Compositionality}: \url{https://arxiv.org/abs/1310.4546} —— 负采样等改进
    \item Pennington et al., 2014. \textit{GloVe: Global Vectors for Word Representation}: \url{https://nlp.stanford.edu/projects/glove/} —— 另一种词向量方法
    \item CS224N 课程官网：\url{https://web.stanford.edu/class/cs224n/}
    \item Jay Alammar, \textit{The Illustrated Word2Vec}: \url{https://jalammar.github.io/illustrated-word2vec/} —— 优秀的可视化教程
\end{itemize}

\end{document}
